\scanpage{57}

% manuscript line(s) 1
\uline{Rev quadratiques en car. $2$}
% manuscript line(s) 2
On a
% manuscript line(s) 2
\[
L\rightrightarrows L^{\otimes2}\qquad\begin{cases}u\longrightarrow u^2\\u\longrightarrow Tu\end{cases}\qquad T\in\Gamma L\ \text{la section trace}
\]
% manuscript line(s) 3
(\uline{NB} ce sont des hom. additifs, mais $u\mapsto u^2$ n'est pas linéaire)
% manuscript line(s) 4
d'où
% manuscript line(s) 4
\[
L\xrightarrow{f}L^{\otimes2}:\quad f(u)=u^2+Tu
\]
% manuscript line(s) 5
et $\check V(L)\underset f\longrightarrow\check V(L^{\otimes2})$ morphisme qui est étale ssi $T$ est inv.
% manuscript line(s) 6
\deleted{On peut considérer} et alors $T:\underline O_X\xrightarrow{\sim}L$, et $f$ peut s'interpréter comme $\underline O_X\longrightarrow\underline O_X:u\longmapsto u^2+u$

% manuscript line(s) 7–8
Considérons \deleted{\unsure{$x'$ et $x''$|$x'$, $x''$}} les deux sections éventuelles $x',x''$ de $X'\subset P_1\subset E$, on a
% manuscript line(s) 9
\[
x''\deleted{\unsure{+}}\unsure{-|+}x'=x''+x'=T\qquad\text{i.e}\qquad x''=x'+T
\]
% manuscript line(s) 10
d'où
% manuscript line(s) 10
\[
{x''}^{(2)}={x'}^{(2)}+T^2\qquad\text{sections de }P_1\mathbin{*}_L L^2\simeq P_1^{(2)}
\]
% manuscript line(s) 11–12
\diagram{\begin{tikzcd}[ampersand replacement=\&,column sep=small]0\arrow[r]\&L\arrow[r]\arrow[d,"{\xi\mapsto\xi^2}"]\&E\arrow[r,"\psi"]\arrow[d,"{x\mapsto x^{(2)}}"]\&\underline O_X\arrow[r]\arrow[d,"{\lambda\mapsto\lambda^2}"]\&0\\0\arrow[r]\&L^{\otimes2}\arrow[r]\&E^{(2)}\arrow[r,"\psi^{(2)}"]\&\underline O_X\arrow[r]\&0\end{tikzcd}}
% manuscript line(s) 12
\margin{$P_1^{(2)}=\psi^{(2)-1}(\{1\})$}
% manuscript line(s) 13
\[
\underset{\mathrm{Sym}^2(E)}{\deleted{\unsure{\mathrm{Sym}^2(E)}}}\big/\left\{(\,x\deleted{\unsure{-}}+y)^{(2)}-\underbrace{x^{(2)}-y^{(2)}}_{xy}\right\}\deleted{\unsure{x^{(2)}=y^{(2)}}}
\]
% manuscript line(s) 14–16
\[
\begin{aligned}\unsure{w|u}\in\Gamma(P_1^{(2)}\wedge T\cdot P_1)&\Longleftrightarrow {\unsure{w|u}}': P_1^{(2)}\simeq TP_1\quad\text{iso de }L^{\otimes2}\text{-torseurs}\\ &\Longleftrightarrow P_1\xrightarrow{u''}P_1^{(2)}\quad\text{compatible avec}\\ &\qquad L\longrightarrow L^{\otimes2}\\ &\qquad u\longmapsto T\cdot u\end{aligned}
\]

% manuscript line(s) 17–20
\uline{NB} $f:\check V(L)\mathbin{-}\check V(L^{\otimes2})$ est fid plat (donc épi fppf) et son noyau est un sous-groupe \inserted{fini} quadratique de $\check V(L)$, soit, \deleted{\unsure{$K$}} $G_T$, étale ssi $T$ est inv. (son ens de ramification est donc $V(T)$, schéma des zéros de $T$)

% manuscript line(s) 21–22
Soit $x\in\deleted{\unsure{\Gamma(P_1)}}\inserted{\Gamma X'}$, i.e. $x\in\Gamma(P_1)$, avec
% manuscript line(s) 22
\[
Q(x)=0\qquad x^2\underset{\underbrace{\text{ou }+}}{-}Tx+N=0
\]
% manuscript line(s) 23
Écrivons $Q(x+u)=0$ \inserted{(où $u\in\Gamma\underline L$)}, \quad $\underset{/}{x^2}+u^2+\underset{/}{Tx}+Tu+\underset{/}{N}=0$, i.e. $u^2+Tu=0$
% archive sheet number: 56
